\section{Overview of Vehicular Data Trading Schemes}\label{sec:lit_data_trading}

{\emergencystretch=2em
\noindent
Recently, vehicular data trading has become a key paradigm for implementing cooperative autonomous driving and smart transportation systems \cite{chen2025lbdt}. The advanced sensors like LiDAR, radar, and high-resolution cameras can produce tremendous amounts of data on local environmental conditions \cite{kang2019toward}. Traditional data trading methods depend on centralized cloud services for data management and pricing. However, centralized systems bring various problems with them, like the possibility of single point of failure, provider lock-in effect, added latency in multihop systems, and extensive chances of tracking user mobility. Decentralized systems (P2P) quickly gained popularity as a scalable option.
}

\section{Blockchain Applications in Vehicular Networks}\label{sec:lit_blockchain}

\noindent
Blockchain technology has provided the means of achieving decentralized and immutable transaction recording and smart contract executing \cite{nakamoto2008bitcoin}. In case of vehicular ad-hoc networks (VANETs), consortium blockchains have been used for ensuring authentication, data sharing, and resource allocation without the use of trusted or centralized authorities \cite{yang2019privacy, firdaus2021blockchain}. For instance, Javaid et al. \cite{javaid2020scalable} proposed a trust management framework based on consortium blockchain for IoV data validation. Michelin et al\cite{michelin2018speedychain}. discussed SpeedyChain, which separates raw payload from blockchain ledger in order to enhance transaction efficiency in smart cities. 
{\emergencystretch=2em
\noindent
Nevertheless, monolithic blockchains do not work in networks with a high density of vehicles because of rapid size of storage, slow propagation of blocks and transaction concluding latencies \cite{liu2018computation}. Dual-chain architecture aimed at making microtransactions site independent from macro-state consensus thus looks very promising \cite{chen2025lbdt}.
}

\section{Reputation Systems and Trust Management}\label{sec:lit_reputation}

\noindent
In the case of vehicular communication, the open nature of the wireless medium facilitates
the transmission of erroneous or fabricated information by malicious or compromised
nodes, affecting the efficiency of functioning of the market\cite{rosenstatter2018modelling}. Traditional models of trust rely on static feedback score or average values, which do not allow them to adapt to the dynamic nature of behavior of nodes, strategies of collusion, and the processes ofwhitewashing \cite{cui2019extensible}. 

Gompertz function has proven to be a good instrument for modeling non-linear and sigmoidal behavior in biological and trust growth processes. By introducing the concept of interaction recency, it becomes possible to make dishonest nodes pay for their malfeasance immediately \cite{chen2025lbdt}. Moreover, the introduction of machine learning classification models, for instance, Logistic Regression or Support Vector Machines allows detecting the instances of non-honest behavior even before they are committed \cite{scikit-learn}.

\section{Double Auction Market Mechanisms}\label{sec:lit_auction}

\noindent
In auction theory a systematic approach is taken to match numerous data sellers with a number of data buyers in cases of dynamic supply and demand conditions \cite{guan2020towards}. when conventional one-sided auctions or fixed prices cannot perform the function of clearing the market when sellers are competing at the same time. 

Double auctions make it possible for buyers to place their bidding price ($BP$) and sellers to place their asking price ($SP$). However, in standard double auctions the assumption is that all data sellers provide data of the same quality and transmission time. In the case of vehicular networks, it becomes much more difficult, as the value of the data falls rapidly in time and the mechanism of the market must include such factors as the reputation penalty and transmission time delay in its calculation of the price for asking \cite{chen2025lbdt}.

\section{Consensus Protocols and Communication Overhead}\label{sec:lit_consensus}

\noindent
Consensus protocols prescribe the method by which decentralized nodes authenticate transactions and finalize block order. Practical Byzantine Fault Tolerance (PBFT) and its variations eliminate the need for energy-intensive PoW (Proof of Work) \cite{lamport1982byzantine}. Nevertheless, standard PBFT has a number of drawbacks, such as $O(N^2)$ message complexity that results in network overload in cases of high vehicular density.
To tackle this issue, a new process has been invented that introduces the usage of a verifiable Random Function (VRF) to randomly select small unpredictable groups of conference participants \cite{xie2023bephap}. If it is utilized along with multi-phase consensus systems like Gosig BFT, the communication difficulties are reduced to $O(N)$ (i.e.,$5N-3$ messages per round)level, allowing confirming the transactions in under 3 seconds even in dense vehicular scenarios \cite{chen2025lbdt}.

\section{Research Gaps Identified}\label{sec:lit_gaps}

\noindent
According to the literature review, four primary research gaps are found in current IoV data trading schemes:
\begin{enumerate}
    \item \textbf{Decoupling Latency from Consensus Complexity:} currently adopted IoV blockchain systems do not achieve sub-3 second confirmation latencies.

    \item \textbf{Static vs. Adaptive Reputation Modeling:}  most of the reputation systems use static thresholds which cannot identify adaptive attack modes that involve shifts between normal and malicious transactions. 

    \item \textbf{Ignoring Latency \& Reputation in Auction Pricing:} current double auction pricing mechanism determines the price based only on monetary bids, without considering data degradation and seller trustworthiness.

    \item \textbf{Lack of Dynamic Sharding for Mobility:} the adoption of fixed geographical zones leads to local bottlenecks, while vehicle clusters move constantly throughout city intersections.
\end{enumerate}

